\chapter{VB\(^{*}\) Main Direction via Hamburger Positivity, Determinacy, and HB\(^{+}\)/BD\(^{+}\)}
\label{v4-arith:w5-j:VBstar-spectral-positivity}

\emph{\Cref{v4-arith:HPAr-vbconverse} isolates the VB\(^{*}\)
converse (zero-placement \(\Rightarrow\) VB\(^{*}\)) as a conditional comparison
between the spectral-regularity hypotheses, operator-positive
Hermite--Biehler de Branges structure, and the B\'aez-Duarte sharp
decay condition. In the opposite direction, the VB\(^{*}\)
main-direction implication says that the four analytic hypotheses
(HB\(^{+}\), BD\(^{+}\), determinacy, positive Hamburger
moment structure) on the primary Fock imply that the non-trivial zeros
of \(\xi_{\mathbb R}\) lie on the critical line. The chain is classical
complex analysis and de Branges space theory: a de Branges datum
whose real part has divisor \(\mathrm{div}\,\Xi\) and whose
Hermite--Biehler kernel is positive forces the zeros of \(\Xi\) to be
real. The positive
Hamburger moment structure identifies the spectral
measure as a positive measure on the real line; Hamburger
determinacy identifies the measure uniquely once it exists; and
HB\(^{+}\) supplies the real-part divisor constraint. Thus
the conclusion is the conditional arrow from these hypotheses to
RH. The hypotheses
HB\(^{+}\), BD\(^{+}\), and determinacy are not consequences of the
Hadamard product; they are the analytic content of the
zero-placement problem.}

\section{Spectral Positivity Context}
\label{v4-arith:w5-j:context}

The operator-theoretic comparison
(Chapter~\ref{v4-arith:HPAr-vbconverse}) reduced the VB\(^{*}\)
converse, under P2 + P3 + H-P\(^{\mathrm{Ar}}\), to the four-way
equivalence
\[
(\mathrm{Det})
\Leftrightarrow
(\mathrm{HB}^{+})
\Leftrightarrow
(\mathrm{BD}^{+})
\Leftrightarrow
\text{zero-placement}\Rightarrow\mathrm{VB}^{*}
\]
of \Cref{v4-arith:hpvbconv:thm:VBstar-converse-spectral-regularity}.
\Cref{v4-arith:hpar-det} reduces the
Hamburger-determinacy component to the sufficient analytic
residual \((\mathrm{R}_{\mathrm{analytic}})\) and obstructed a tempting
Bombieri--Lagarias shortcut. Neither argument addresses the
main-direction implication
\[
(\mathrm{HB}^{+})+(\mathrm{BD}^{+})+(\mathrm{Det})
\;\Longrightarrow\;
\text{spectral zeros of \(\xi_{\mathbb R}\) on the critical line},
\]
which is the sufficient direction of the analytic hypothesis of
\Cref{v4-arith:platonic:def:RH-package}. The proof is classical
complex analysis and
de Branges space theory in the sense of the 1968 monograph;
no arithmetic input beyond the Hadamard factorisation of
\(\xi_{\mathbb R}\) is used.

\subsection{Domain}
\label{v4-arith:w5-j:domain}

The theorem is the implication from
(HB\(^{+}\)) + (BD\(^{+}\)) + (Det) + positive Hamburger moment
structure to spectral zero-placement on the critical line.

The analytic inputs are HB\(^{+}\), BD\(^{+}\), (Det), and
positive Hamburger moment structure on \(\xi_{\mathbb R}\).
These are RH-equivalent in the Hermite--Biehler case and spectral
regularity inputs in the moment-problem cases. They are the subject of
Chapters~\ref{v4-arith:HPAr-vbconverse}
and~\ref{v4-arith:hpar-det}.

\subsection{Independent Sources}
\label{v4-arith:w5-j:sources}

The source list for the VB\(^{*}\) converse uses
de Branges 1968, Lagarias 1999, B\'aez-Duarte 2003, Burnol 2000,
Akhiezer 1965, Simon 1998, Shohat--Tamarkin 1943.  The
main-direction implication inherits this list and adds the
classical Hamburger existence theorem, Lagarias' 1999 de Branges
survey, Conrey--Li 2000, and Beurling--Malliavin density. The
proof sources of constructions 2 and 3 (determinacy reductions,
operator-positivity reformulations) are kept disjoint from the
complex-analysis sources below:

\begin{description}
\item[de Branges, \emph{Hilbert Spaces of Entire Functions}
(Prentice-Hall 1968).] Foundational monograph; theorems 22--23
(Hermite--Biehler structure), theorem 32 (multiplication operator
on \(\mathcal{H}(E)\)), theorems 40--41 (invariant subspaces of
the shift).
\item[Lagarias, Acta Arith.~89 (1999).] Survey of de Branges
reformulations of RH; explicit statement of the main-direction
theorem for entire functions of order 1.
\item[Conrey--Li, arXiv:math/9812166 (2000).] Construction of
de Branges spaces from \(L\)-functions; critical discussion of the
Hermite--Biehler condition on \(E_{\xi}\); counter-examples in the
non-self-adjoint case.
\item[Beurling--Malliavin, Acta Math.~107 (1962) 291--309.]
Density theorem: the completeness radius of an exponential system
\(\{e^{i\lambda_{k}x}\}\) is computed by Beurling--Malliavin
density of \(\{\lambda_{k}\}\).
\item[Hamburger, Math.~Ann.~81 (1920) 235--319.] Original
existence theorem: a Hamburger moment problem has a positive
representing measure iff all Hankel matrices are
positive-semidefinite.
\item[Levin, \emph{Distribution of Zeros of Entire Functions}
(AMS 1980), Ch.~V--VII.] Classical reference on Hermite--Biehler,
Hermite--Kakeya--Obreshkov; entire functions of exponential type.
\end{description}

Independence of inputs: the complex analysis (Hadamard
factorisation of \(\xi_{\mathbb R}\), de Branges structure theorem,
Beurling--Malliavin density) is not used in the
\(\mathrm{VB}^{*}\)-converse comparison; and that comparison
(operator-positivity
reformulations, Jacobi-matrix spectral theory, B\'aez-Duarte decay)
is not used in the main-direction implication. The two lists share two foundational references: de Branges 1968
(Hermite--Biehler theory) and Akhiezer 1965
(moment theory, used for the Jacobi matrix and for Hamburger
existence and Carleman's criterion). These shared
references serve logically distinct roles: Akhiezer supplies
operator spectral methods on one side and classical moment
existence.
Simon 1998 (spectral theory of Jacobi matrices) appears in the
operator-theoretic list and is not cited below; Shohat--Tamarkin 1943
is likewise absent below. The two argument paths are
therefore independent on the operative arguments, with the
shared foundational references serving independent functions in the
two constructions.

\section{Primary Fock and \(\xi_{\mathbb R}\) as Character}
\label{v4-arith:w5-j:primary-Fock}

Fix the primary Fock
\(\mathcal{V}^{\mathrm{prim}}_{\mathrm{Heis},\mathrm{prim}}\) of
Chapters~\ref{v4-arith:rh-reformulation}
and~\ref{v4-arith:HPAr-vbconverse}: the primary sector of the
arithmetic Heisenberg Fock, equipped with the Petersson inner
product of~\Cref{v4-arith:kp:cor:P3-unconditional}. Its character
on the archimedean Tate fibre is the completed zeta function
\[
\Lambda_{\zeta}(s)
\;=\;\pi^{-s/2}\,\Gamma(s/2)\,\zeta(s),
\]
meromorphic with simple poles at \(s=0\) and \(s=1\); the entire
Li--Hadamard form is
\[
\xi_{\mathbb R}(s)\;:=\;\tfrac12\,s(s-1)\,\Lambda_{\zeta}(s).
\]
The factor \(s(s-1)\) removes the Tate poles; the Gamma factor
cancels the negative even trivial zeros of \(\zeta\); the remaining
zero set is exactly the non-trivial zero set of \(\zeta\). This is
the convention fixed in
\Cref{v4-arith:rh:rem:equivalence-to-RH} and
\Cref{v4-arith:rh:rem:trivial-zero}.

\begin{proposition}[\(\xi_{\mathbb R}\) is entire of order \(1\), Hadamard genus \(1\)]
\label{v4-arith:w5-j:prop:xi-order}
\ClaimStatusProvedElsewhere (Edwards, \emph{Riemann's Zeta Function}
(Academic Press 1974), \S1.9). The function
\(\xi_{\mathbb R}(s)\) is entire of order \(1\) and Hadamard genus \(1\):
\begin{equation}
\label{v4-arith:w5-j:eq:hadamard}
\xi_{\mathbb R}(s)
\;=\;e^{a+bs}\prod_{\rho}\Bigl(1-\frac{s}{\rho}\Bigr)e^{s/\rho},
\end{equation}
with explicit constants
\(a=-\log 2\),
\(b=-1-\tfrac{1}{2}\gamma_{\mathrm{E}}+\tfrac{1}{2}\log(4\pi)\),
product over non-trivial zeros \(\rho\), and
\(\sum_{\rho}|\rho|^{-2}<+\infty\).
\end{proposition}

\begin{proof}
Classical: the functional equation and
\(|\xi_{\mathbb R}(s)|\leq\exp(C|s|\log|s|)\) uniformly in vertical
strips (Stirling bounds on \(\Gamma\), polynomial control on
\(\zeta\) away from the pole) give order \(1\). Hadamard genus \(1\)
follows from the zero density estimate of Riemann--von Mangoldt,
which gives \(\sum_{\rho}|\rho|^{-1-\epsilon}<+\infty\) for every
\(\epsilon>0\) but
\(\sum_{\rho}|\rho|^{-1}=+\infty\). The Hadamard factorisation
theorem then gives~\eqref{v4-arith:w5-j:eq:hadamard} with genus \(1\).
At \(s=0\), the product is \(1\) and
\(\xi_{\mathbb R}(0)=1/2\), so \(a=-\log2\).  Differentiating
\(\log\xi_{\mathbb R}(s)\) at \(s=0\), using
\(\zeta(0)=-1/2\), \(\zeta'(0)=-(1/2)\log(2\pi)\), and
\(\Gamma(s/2)=2s^{-1}-\gamma_{\mathrm E}+O(s)\), gives
\[
\frac{\xi_{\mathbb R}'(0)}{\xi_{\mathbb R}(0)}
=-1-\frac{\gamma_{\mathrm E}}2+\frac12\log(4\pi),
\]
which is \(b\).
\end{proof}

\begin{proposition}[functional equation]
\label{v4-arith:w5-j:prop:func-eq}
\ClaimStatusProvedElsewhere. The entire function
\(\xi_{\mathbb R}\) satisfies the functional equation
\begin{equation}
\label{v4-arith:w5-j:eq:func-eq}
\xi_{\mathbb R}(s)\;=\;\xi_{\mathbb R}(1-s),
\qquad s\in\mathbb{C}.
\end{equation}
Equivalently, \(\xi_{\mathbb R}(\tfrac{1}{2}+iz)\) is an entire
function of \(z\) with real Taylor coefficients.
\end{proposition}

\begin{proof}
The functional equation for \(\Lambda_{\zeta}\) is Riemann 1859 /
Tate 1950. The extension to \(\xi_{\mathbb R}\) follows because
\(s(s-1)=(1-s)(-s)=(1-s)((1-s)-1)\) is invariant under
\(s\mapsto 1-s\). The substitution \(s=\tfrac{1}{2}+iz\) turns
\(s\mapsto 1-s\) into \(z\mapsto -z\); combined with
\(\xi_{\mathbb R}(\overline{s})=\overline{\xi_{\mathbb R}(s)}\)
(real on the real axis), this gives
\(\xi_{\mathbb R}(\tfrac{1}{2}+iz)=\overline{\xi_{\mathbb R}(\tfrac{1}{2}-i\overline{z})}=\overline{\xi_{\mathbb R}(\tfrac{1}{2}+iz)}\)
for \(z\in\mathbb{R}\); hence the Taylor coefficients at \(z=0\) are real.
\end{proof}

\begin{definition}[shifted character]
\label{v4-arith:w5-j:def:shifted-character}
\ClaimStatusDefinitional. Write
\[
\Xi(z)\;:=\;\xi_{\mathbb R}\!\left(\tfrac{1}{2}+iz\right),
\qquad z\in\mathbb{C}.
\]
Then \(\Xi\) is entire of order \(1\) with real Taylor coefficients.
Its zero set is
\[
\{\gamma_{\rho}:\rho=\tfrac{1}{2}+i\gamma_{\rho}\text{ non-trivial zero of }\xi_{\mathbb R}\},
\]
symmetric under \(z\mapsto-z\) by~\eqref{v4-arith:w5-j:eq:func-eq}.
Under RH the zeros \(\gamma_{\rho}\) are real; unconditionally they
are complex and the \(z\mapsto-z\) symmetry becomes
\(\gamma\leftrightarrow-\gamma\) as a multi-set.
\end{definition}

\begin{remark}[the orientation of the problem]
\label{v4-arith:w5-j:rem:orientation}
RH is the statement that the zero set of \(\Xi\) is a subset of the
real line. The Hadamard and functional-equation facts
(\Cref{v4-arith:w5-j:prop:xi-order},
\Cref{v4-arith:w5-j:prop:func-eq}) are unconditional; they
constrain the zero set to be symmetric under \(z\mapsto-z\) but not
to be real. The role of the analytic hypotheses (HB\(^{+}\),
BD\(^{+}\), determinacy) is to force this realness by real-part
divisor control and spectral-measure rigidity.
\end{remark}

\section{Hadamard Factorisation and the Zero Measure}
\label{v4-arith:w5-j:zero-measure}

\begin{definition}[primary zero multi-set]
\label{v4-arith:w5-j:def:primary-zero-set}
\ClaimStatusDefinitional. The \emph{primary zero multi-set} is
\[
\mathcal{Z}_{\xi}
\;:=\;\{\gamma_{\rho}\in\mathbb{C}:\rho=\tfrac{1}{2}+i\gamma_{\rho}\text{ non-trivial zero of }\xi_{\mathbb R}\},
\]
with multiplicity. Under the symmetrisation
\(\rho\leftrightarrow 1-\rho\), the set \(\mathcal{Z}_{\xi}\) is closed
under \(\gamma\leftrightarrow-\gamma\). Under the complex conjugation
\(\rho\leftrightarrow\overline{\rho}\), it is closed under
\(\gamma\leftrightarrow-\overline{\gamma}\); combining, it is closed
under \(\gamma\leftrightarrow\overline{\gamma}\) and symmetric under the Klein four group
\(\{1,\gamma\mapsto-\gamma,\gamma\mapsto\overline{\gamma},\gamma\mapsto-\overline{\gamma}\}\)
on the whole plane.
\end{definition}

\begin{definition}[primary zero counting measure]
\label{v4-arith:w5-j:def:zero-counting-measure}
\ClaimStatusDefinitional. The \emph{primary zero counting measure}
is the multi-set measure
\[
\mu_{\xi}\;:=\;\sum_{\gamma\in\mathcal{Z}_{\xi}}\delta_{\gamma}
\]
on \(\mathbb{C}\) (with multiplicity). Under RH, \(\mu_{\xi}\) is
supported on \(\mathbb{R}\). Under the functional
equation~\eqref{v4-arith:w5-j:eq:func-eq}, \(\mu_{\xi}\) is
\(z\mapsto-z\)-invariant. Under the reality of Taylor coefficients
(\Cref{v4-arith:w5-j:prop:func-eq}), \(\mu_{\xi}\) is
\(z\mapsto\overline{z}\)-invariant. The invariant group is thus the
Klein four group on \(\mathbb{C}\).
\end{definition}

\begin{definition}[regularised moment sequence]
\label{v4-arith:w5-j:def:moment-sequence}
\ClaimStatusDefinitional. For \(n\geq 0\), the \emph{regularised
\(n\)-th moment of \(\mu_{\xi}\)} is
\begin{equation}
\label{v4-arith:w5-j:eq:moment-def}
M_{n}\;:=\;\sum_{\gamma\in\mathcal{Z}_{\xi}}^{\mathrm{reg}}\gamma^{n}
\end{equation}
in the sense of Bost--Connes regularisation (zeta-regularised
summation in the cutoff height \(T\), with the arithmetic Euler
factor subtracted). The regularisation exists and is finite for
every \(n\geq 0\) by the Riemann--von Mangoldt counting formula
\(N(T)\sim(T/2\pi)\log(T/2\pi)\) and the
Hadamard-factorisation estimate
\(\sum_{\gamma}|\gamma|^{-1-\epsilon}<+\infty\) for every
\(\epsilon>0\).
\end{definition}

\begin{proposition}[reality of moments]
\label{v4-arith:w5-j:prop:moments-real}
\ClaimStatusProvedHere. For every \(n\geq 0\), \(M_{n}\in\mathbb{R}\).
The odd moments \(M_{2k+1}\) vanish under the \(\gamma\mapsto-\gamma\)
symmetry.
\end{proposition}

\begin{proof}
\(\mu_{\xi}\) is \(\gamma\mapsto\overline{\gamma}\)-invariant, so for
every \(n\) the complex conjugate of
\(\sum_{\gamma}\gamma^{n}\) equals
\(\sum_{\gamma}\overline{\gamma}^{n}=\sum_{\gamma}\gamma^{n}\)
(matching the conjugate partner within the multi-set). Hence
\(M_{n}\in\mathbb{R}\). For odd \(n=2k+1\), the \(\gamma\mapsto-\gamma\)
symmetry pairs each \(\gamma^{2k+1}\) with \((-\gamma)^{2k+1}=-\gamma^{2k+1}\),
giving cancellation; the regularisation preserves this cancellation.
\end{proof}

\begin{remark}[positivity of the measure is strictly weaker than RH]
\label{v4-arith:w5-j:rem:positivity-weaker-than-RH}
The measure \(\mu_{\xi}\) as a multi-set measure on \(\mathbb{C}\) is
always positive (it is a counting measure). The question is
whether its restriction to \(\mathbb{R}\) carries the full mass; that
is the content of RH. Positivity of \(\mu_{\xi}\) as a measure on
\(\mathbb{C}\) gives no information about whether the support is on
\(\mathbb{R}\). This is the reason a positive Hamburger
moment representation (a positive representation on \(\mathbb{R}\)
with prescribed real moments) is strictly finer than mere
positivity of \(\mu_{\xi}\).
\end{remark}

\section{Hamburger Moment Problem for \(\mu_{\xi}\)}
\label{v4-arith:w5-j:SH-moment}

\begin{definition}[Hamburger moment problem for \(\mu_{\xi}\)]
\label{v4-arith:w5-j:def:hamburger-problem}
\ClaimStatusDefinitional. The \emph{Hamburger moment problem} for
the sequence \(\{M_{n}\}_{n\geq 0}\) of
\Cref{v4-arith:w5-j:def:moment-sequence} is: find a positive Borel
measure \(\nu_{\xi}\) on \(\mathbb{R}\) with
\begin{equation}
\label{v4-arith:w5-j:eq:hamburger-representation}
M_{n}\;=\;\int_{\mathbb{R}}x^{n}\,d\nu_{\xi}(x)\qquad\forall n\geq 0.
\end{equation}
The problem is \emph{solvable} if such \(\nu_{\xi}\) exists, and
\emph{positive Hamburger} if \(\nu_{\xi}\) is a positive measure on
\(\mathbb R\). The Stieltjes problem belongs to the square-pushforward
\(x\mapsto x^{2}\), not to the signed ordinate moments themselves.
\end{definition}

\begin{theorem}[Hamburger's existence theorem]
\label{v4-arith:w5-j:thm:hamburger-existence}
\ClaimStatusProvedElsewhere (Hamburger, Math.~Ann.~81 (1920);
Akhiezer 1965 \S2.1). A moment problem
\eqref{v4-arith:w5-j:eq:hamburger-representation} is solvable if
and only if the infinite Hankel matrix
\[
H_{\infty}\;=\;(M_{i+j})_{i,j\geq 0}
\]
is positive-semidefinite:
\(\sum_{i,j}c_{i}\overline{c_{j}}M_{i+j}\geq 0\) for every finitely
supported sequence \((c_{i})\). Equivalently, every finite truncation
\(H_{N}=(M_{i+j})_{0\leq i,j\leq N}\) has non-negative principal
minors.
\end{theorem}

\begin{definition}[HB\(^{+}\)/BD\(^{+}\)/Det hypothesis]
\label{v4-arith:w5-j:def:spectral-package}
\ClaimStatusDefinitional. The \emph{spectral positivity hypothesis}
on the primary Fock is the conjunction of:
\begin{itemize}
\item[(SH)] \emph{Positive Hamburger moment sequence.}
The Hankel matrix \(H_{\infty}=(M_{i+j})_{i,j\geq 0}\) of
\Cref{v4-arith:w5-j:thm:hamburger-existence} is
positive-semidefinite.
\item[(Det)] \emph{Hamburger determinacy.}
\(\sum_{n\geq 1}M_{2n}^{-1/(2n)}=+\infty\) (Carleman's criterion;
sufficient for determinacy of \(\nu_{\xi}\)).
\item[(HB\(^{+}\))] \emph{Hermite--Biehler positivity with divisor
control.} There is a de Branges datum
\(E_{\xi}(z)=A_{\xi}(z)-iB_{\xi}(z)\), with \(A_{\xi},B_{\xi}\)
real entire, such that \(E_{\xi}\) is Hermite--Biehler, the
associated kernel \(K_{E_{\xi}}\) is positive-definite, and the
real part \(A_{\xi}=(E_{\xi}+E_{\xi}^{*})/2\) has divisor
\(\mathrm{div}\,\Xi\), with multiplicity.
\item[(BD\(^{+}\))] \emph{Beurling--Malliavin super-critical density.}
After the HB\(^{+}\) real-part theorem has made
\(\mathcal Z_{\xi}\) real, the associated exponential system
\(\{e^{i\gamma x}\}_{\gamma\in\mathcal Z_{\xi}}\) is complete in
\(L^{2}((-\pi,\pi))\). A sufficient source is
\(D_{\mathrm{BM}}(\mathcal Z_{\xi})>1\); at equality the endpoint
completeness assertion is separate data.
\end{itemize}
\end{definition}

\begin{proposition}[(SH) implies real-positive representing measure]
\label{v4-arith:w5-j:prop:SH-positive-measure}
\ClaimStatusProvedHere. Under (SH), the moment problem
\eqref{v4-arith:w5-j:eq:hamburger-representation} admits at least
one positive Borel measure \(\nu_{\xi}\) on \(\mathbb{R}\) solving it.
\end{proposition}

\begin{proof}
Direct application of Hamburger's existence theorem
(\Cref{v4-arith:w5-j:thm:hamburger-existence}) to the moment
sequence \(\{M_{n}\}_{n\geq 0}\) under the positive-semidefiniteness
hypothesis (SH).
\end{proof}

\begin{remark}[positivity is a genuine arithmetic input]
\label{v4-arith:w5-j:rem:SH-arithmetic-input}
The hypothesis (SH) is not a tautology. The moments \(M_{n}\) are
real-valued regularised sums of complex \(n\)-th powers
(\Cref{v4-arith:w5-j:def:moment-sequence}); positive-semidefiniteness
of the Hankel matrix is what transfers the complex multi-set
\(\mathcal{Z}_{\xi}\) to a real positive measure. This is the
content Hamburger existence requires; it is distinct from the
question whether the zeros are real.
\end{remark}

\section{Hamburger Determinacy and Carleman}
\label{v4-arith:w5-j:carleman-det}

\begin{theorem}[Carleman's determinacy criterion]
\label{v4-arith:w5-j:thm:carleman}
\ClaimStatusProvedElsewhere (Carleman 1926; Akhiezer 1965 \S2.4).
Let \(\{M_{n}\}_{n\geq 0}\) be a Hamburger moment sequence with
positive-semidefinite Hankel matrix and \(M_{2n}>0\) for all \(n\geq
0\). If
\begin{equation}
\label{v4-arith:w5-j:eq:carleman-sum}
\sum_{n\geq 1}M_{2n}^{-1/(2n)}\;=\;+\infty,
\end{equation}
then the representing measure \(\nu_{\xi}\) is unique
(\emph{determinate}).
\end{theorem}

\begin{proposition}[Carleman is not supplied by zero density]
\label{v4-arith:w5-j:prop:carleman-RH-heuristic}
\ClaimStatusProvedHere. The Riemann--von Mangoldt zero-counting law
and the Bombieri--Lagarias--Conrey estimates do not imply
\eqref{v4-arith:w5-j:eq:carleman-sum}. Their available consequence is
only the Bertrand-scale comparison
\[
M_{2n}^{-1/(2n)}
\;\gtrsim\;
\frac{c}{n(\log n)^2},
\qquad n\geq 2,
\]
which is summable and therefore cannot force Carleman divergence.
\end{proposition}

\begin{proof}
The computation is the one recorded in
\Cref{v4-arith:r-comparison:claim-3}. The logarithmic inflation in the
regularised moment estimate gives
\((\log n)^{-2}\) after the \(1/(2n)\)-root and inversion. Thus the
comparison series is \(\sum_{n\geq2}1/(n(\log n)^2)\), which
converges. A summable lower comparison is silent about the divergence
of the Carleman sum.
\end{proof}

\begin{remark}[the heuristic is not a theorem]
\label{v4-arith:w5-j:rem:heuristic-not-theorem}
The moment asymptotics \(M_{2n}\asymp C^{n}n!^{2}\) are a
saddle-point heuristic. If they were replaced by a rigorous two-sided
bound for a positive moment functional, Carleman's series would have
harmonic size. Turning that statement into a theorem is the content of
the residual
\((\mathrm{R}_{\mathrm{analytic}})\) of
\Cref{v4-arith:hpar-det:thm:single-analytic-residual}. The
main-direction implication uses (Det) as a \emph{hypothesis}: the Carleman sum
is assumed to diverge. Verifying this divergence is a separate
analytic problem.
\end{remark}

\begin{corollary}[under (SH) + (Det) the measure \(\nu_{\xi}\) is unique]
\label{v4-arith:w5-j:cor:unique-measure}
\ClaimStatusProvedHere. Under (SH) + (Det) of
\Cref{v4-arith:w5-j:def:spectral-package}, there exists a unique
positive Borel measure \(\nu_{\xi}\) on \(\mathbb{R}\) with moments
\eqref{v4-arith:w5-j:eq:hamburger-representation}.
\end{corollary}

\begin{proof}
(SH) + Hamburger existence
(\Cref{v4-arith:w5-j:thm:hamburger-existence}) gives solvability.
Carleman (\Cref{v4-arith:w5-j:thm:carleman}) under (Det) gives
uniqueness. The conjunction is the claim.
\end{proof}

\section{Hermite--Biehler and the de Branges Space \(\mathcal{H}(E_{\xi})\)}
\label{v4-arith:w5-j:HB-deBranges}

\begin{definition}[de Branges datum over \(\Xi\)]
\label{v4-arith:w5-j:def:HB-decomposition}
\ClaimStatusDefinitional. A \emph{de Branges datum over \(\Xi\)} is a
triple \((E_{\xi},A_{\xi},B_{\xi})\) with
\begin{equation}
\label{v4-arith:w5-j:eq:HB-decomp}
E_{\xi}(z)\;=\;A_{\xi}(z)-iB_{\xi}(z),
\end{equation}
where \(A_{\xi}\) and \(B_{\xi}\) are real entire functions of finite
exponential type.  The datum is required to carry a real-part divisor equality
\[
\mathrm{div}(A_{\xi})\;=\;\mathrm{div}(\Xi)
\]
on \(\mathbb{C}\), with multiplicities.  It is not obtained by
setting \(E_{\xi}=\Xi\); for the real entire function \(\Xi\) one has
\(\Xi^{*}=\Xi\), so the strict Hermite--Biehler inequality would be
impossible.
\end{definition}

\begin{remark}[conventions match de Branges 1968]
\label{v4-arith:w5-j:rem:de-branges-convention}
The notation \(E=A-iB\) with \(A\) and \(B\) real-on-real is de
Branges' standard form for candidate Hermite--Biehler functions. The
condition for Hermite--Biehler is
\(|E(z)|>|E^{*}(z)|\) for \(\mathrm{Im}\,z>0\), where
\(E^{*}(z)=\overline{E(\overline{z})}\).  The divisor equality for \(A_{\xi}\)
is part of the arithmetic datum; Hermite--Biehler positivity of an
unrelated entire function says nothing about zeros of \(\Xi\).
\end{remark}

\begin{definition}[de Branges space \(\mathcal{H}(E_{\xi})\)]
\label{v4-arith:w5-j:def:dB-space}
\ClaimStatusDefinitional. Assuming \(E_{\xi}=A_{\xi}-iB_{\xi}\) is
Hermite--Biehler (see (HB\(^{+}\)) below), the
\emph{de Branges space} \(\mathcal{H}(E_{\xi})\) is the
reproducing-kernel Hilbert space of entire functions \(f\) with
\(f/E_{\xi},f^{*}/E_{\xi}\in H^{2}(\mathbb{C}_{+})\), with
reproducing kernel
\begin{equation}
\label{v4-arith:w5-j:eq:dB-kernel}
K_{E_{\xi}}(w,z)
\;=\;\frac{\overline{E_{\xi}(w)}E_{\xi}(z)-\overline{E_{\xi}^{*}(w)}E_{\xi}^{*}(z)}{2\pi i(\overline{w}-z)}.
\end{equation}
\end{definition}

\begin{proposition}[Hermite--Biehler kernel positivity]
\label{v4-arith:w5-j:prop:HBplus-iff-kernel-positive}
\ClaimStatusProvedElsewhere (de Branges 1968 Thm.~22;
Lagarias 1999 Acta Arith.~89 \S3). For an entire function
\(E=A-iB\) of bounded type in \(\mathbb C_{+}\), the following are
equivalent:
\begin{enumerate}
\item \(E\) is Hermite--Biehler: \(|E(z)|>|E^{*}(z)|\)
for \(\mathrm{Im}\,z>0\).
\item The reproducing kernel \(K_{E}(w,z)\) of
\eqref{v4-arith:w5-j:eq:dB-kernel} is positive-semidefinite:
\(\sum_{j,k}c_{j}\overline{c_{k}}K_{E}(w_{k},w_{j})\geq 0\)
for every finite collection \((w_{j},c_{j})\).
\end{enumerate}
If these conditions hold, the zeros of the real entire functions
\(A=(E+E^{*})/2\) and \(B=(E-E^{*})/(2i)\) are real; if \(E\) has no
real zeros, these zeros are simple and interlace.
\end{proposition}

\begin{proof}
(1) \(\Leftrightarrow\) (2): de Branges 1968 Thm.~22 proves that
\(K_{E}\) is positive-semidefinite exactly when \(E\) is
Hermite--Biehler. The forward implication uses the
Szeg\H{o}/Blaschke factorisation of \(|E|^{2}-|E^{*}|^{2}\) as a
positive harmonic function on \(\mathbb{C}_{+}\); the converse uses
that positive-semidefiniteness of \(K_{E}\) at the specific points
\(w=\overline{z}\) implies \(|E(z)|^{2}\geq|E^{*}(z)|^{2}\) strictly
in the open upper half plane.
The last assertion is the Hermite--Biehler separation theorem.
\end{proof}

\begin{theorem}[Hermite--Biehler real-part control forces zero placement]
\label{v4-arith:w5-j:thm:zeros-of-Xi-under-HB}
\ClaimStatusProvedHere. Under (HB\(^{+}\)), the zero multi-set
\(\mathcal{Z}_{\xi}\) of \(\Xi\) is contained in \(\mathbb R\).
\end{theorem}

\begin{proof}
By (HB\(^{+}\)), \(E_{\xi}=A_{\xi}-iB_{\xi}\) is
Hermite--Biehler and \(\mathrm{div}(A_{\xi})=\mathrm{div}(\Xi)\).
The Hermite--Biehler separation theorem gives that every zero of
\(A_{\xi}\) is real. Hence every zero of \(\Xi\) is real.
\end{proof}

\begin{corollary}[zero multi-set is real under HB\(^{+}\)]
\label{v4-arith:w5-j:cor:no-UHP-zeros}
\ClaimStatusProvedHere. Under (HB\(^{+}\)), no non-trivial zero of
\(\xi_{\mathbb R}\) has \(\mathrm{Im}\,\gamma\neq0\) in the
coordinate \(\rho=\tfrac12+i\gamma\).
\end{corollary}

\begin{proof}
\Cref{v4-arith:w5-j:thm:zeros-of-Xi-under-HB} gives
\(\gamma\in\mathbb R\) for every zero.
\end{proof}

\begin{remark}[HB\(^{+}\) is an RH-strength datum]
\label{v4-arith:w5-j:rem:HBplus-alone}
\Cref{v4-arith:w5-j:cor:no-UHP-zeros} shows that (HB\(^{+}\)),
including real-part divisor control over \(\Xi\), forces every non-trivial zero
\(\gamma\) of \(\Xi\) to be real. Translated back via
\(\rho=\tfrac{1}{2}+i\gamma\), this is exactly
\(\mathrm{Re}\,\rho=\tfrac{1}{2}\).  Thus (HB\(^{+}\)) is not a
soft positivity assumption: it already contains the de Branges
real-part mechanism at RH strength.
\end{remark}

\section{BD\(^{+}\): Beurling--Malliavin Super-Critical Density}
\label{v4-arith:w5-j:BDplus}

\begin{definition}[Beurling--Malliavin sampling density]
\label{v4-arith:w5-j:def:BM-density}
\ClaimStatusDefinitional (Beurling--Malliavin, Acta Math.~107
(1962)). For a real sequence \(\Lambda=\{\lambda_{k}\}\subset\mathbb{R}\)
with \(|\lambda_{k}|\to\infty\), the
\emph{Beurling--Malliavin sampling density}
\(D_{\mathrm{BM}}(\Lambda)\) is the infimum of \(a>0\) for which there
exists a real measurable function \(\phi:\mathbb{R}\to\mathbb{R}\)
with
\[
\int_{\mathbb{R}}\frac{|\phi(x)|}{1+x^{2}}\,dx<+\infty,
\qquad
\phi(\lambda_{k})\geq a|\lambda_{k}|-O(1)
\quad\text{for large }|\lambda_{k}|.
\]
For complex \(\gamma\in\mathbb{C}\), define
\(D_{\mathrm{BM}}(\{\gamma_{k}\})\) by projecting to the real part
and applying the real-line definition.
\end{definition}

\begin{theorem}[Beurling--Malliavin completeness theorem]
\label{v4-arith:w5-j:thm:BM-completeness}
\ClaimStatusProvedElsewhere (Beurling--Malliavin, Acta Math.~107
(1962); Koosis, \emph{The Logarithmic Integral II} (Cambridge 1992)
Ch.~VIII). Let \(\Lambda=\{\lambda_{k}\}\subset\mathbb{R}\). The
exponential system \(\{e^{i\lambda_{k}x}\}\) is complete in
\(L^{2}((-\pi a,\pi a))\) for every \(a<D_{\mathrm{BM}}(\Lambda)\) and
incomplete for every \(a>D_{\mathrm{BM}}(\Lambda)\). At the critical
value \(a=D_{\mathrm{BM}}(\Lambda)\) the completeness question is
subtle.
\end{theorem}

\begin{proposition}[BD\(^{+}\) gives uniqueness sampling]
\label{v4-arith:w5-j:prop:BDplus-reconstruction}
\ClaimStatusProvedHere. Under (BD\(^{+}\))
and (HB\(^{+}\))
(\Cref{v4-arith:w5-j:def:spectral-package}), the evaluation map
\begin{equation}
\label{v4-arith:w5-j:eq:sampling-map}
\mathcal{H}(E_{\xi})\;\to\;\ell^{2}(\mathcal{Z}_{\xi};|K_{E_{\xi}}|^{-1}),
\qquad
f\mapsto(f(\gamma))_{\gamma\in\mathcal{Z}_{\xi}},
\end{equation}
is injective on the Paley--Wiener realisation of
\(\mathcal H(E_{\xi})\). If the stronger sampling inequality is
included, the values on \(\mathcal Z_{\xi}\) determine \(f\) by the
corresponding de Branges interpolation formula.
\end{proposition}

\begin{proof}
The exponential type of \(E_{\xi}\) is \(\pi\) (normalised so that
\(\Xi\) has exponential type \(\pi\); Lagarias 1999 \S4.2).
Beurling--Malliavin
(\Cref{v4-arith:w5-j:thm:BM-completeness}) gives completeness for
strict density and leaves the endpoint to an explicit completeness
assumption. This is exactly the BD\(^{+}\) clause. Under (HB\(^{+}\))
the set \(\mathcal Z_{\xi}\) is real, so the exponentials are defined
on \(L^{2}((-\pi,\pi))\). The de Branges space
\(\mathcal{H}(E_{\xi})\) is the Paley--Wiener realisation of this
space under the Mellin--Fourier transform associated with \(E_{\xi}\)
(de Branges 1968 Ch.~III; Lagarias 1999 \S4). Completeness says that
no non-zero vector is orthogonal to all exponentials, which is
injectivity of the sampling map. A reconstruction formula requires the
stronger sampling inequality, hence is part of the strengthened input.
\end{proof}

\begin{corollary}[zero set is sampling for \(\mathcal{H}(E_{\xi})\)]
\label{v4-arith:w5-j:cor:zeros-sample}
\ClaimStatusProvedHere. Under (HB\(^{+}\)) + (BD\(^{+}\)), the zero
multi-set satisfies \(\mathcal{Z}_{\xi}\subset\mathbb{R}\) and is a
uniqueness family for \(\mathcal{H}(E_{\xi})\): no non-zero element
of \(\mathcal{H}(E_{\xi})\) vanishes on all of \(\mathcal{Z}_{\xi}\).
If the strengthened sampling inequality is imposed, the values on
\(\mathcal Z_{\xi}\) reconstruct the function.
\end{corollary}

\begin{proof}
(HB\(^{+}\)) and \Cref{v4-arith:w5-j:cor:no-UHP-zeros} give
\(\mathcal{Z}_{\xi}\subset\mathbb{R}\). (BD\(^{+}\)) and
\Cref{v4-arith:w5-j:prop:BDplus-reconstruction} give injectivity of
the evaluation map.
\end{proof}

\section{Main Theorem: (HB\(^{+}\)) + (BD\(^{+}\)) + (Det) \(\Rightarrow\) RH}
\label{v4-arith:w5-j:main-theorem}

\begin{theorem}[VB\(^{*}\) main-direction implication]
\label{v4-arith:w5-j:thm:VBstar-main-direction}
\ClaimStatusProvedHere. Assume the spectral positivity hypothesis
(SH) + (Det) + (HB\(^{+}\)) + (BD\(^{+}\)) of
\Cref{v4-arith:w5-j:def:spectral-package} on the primary Fock.
Then all non-trivial zeros of \(\xi_{\mathbb R}\) lie on the
critical line \(\mathrm{Re}\,s=\tfrac{1}{2}\). Equivalently,
\(\mathcal{Z}_{\xi}\subset\mathbb{R}\); equivalently, RH holds.
\end{theorem}

\begin{proof}
By (HB\(^{+}\)) and
\Cref{v4-arith:w5-j:thm:zeros-of-Xi-under-HB},
\(\mathcal{Z}_{\xi}\subset\mathbb{R}\).
The definition of \(\mathcal{Z}_{\xi}\) then gives
\(\rho=\tfrac{1}{2}+i\gamma_{\rho}\) with
\(\gamma_{\rho}\in\mathbb{R}\) for every non-trivial zero of
\(\xi_{\mathbb R}\).  Thus \(\mathrm{Re}\,\rho=\tfrac{1}{2}\).
\end{proof}

\begin{remark}[the logical order of the hypotheses]
\label{v4-arith:w5-j:rem:proof-order}
The zero-placement implication uses the real-part divisor control in
(HB\(^{+}\)).  The clauses (SH), (Det),
and (BD\(^{+}\)) belong to the stronger spectral-measure hypothesis:
they concern positivity, uniqueness, and sampling of the zero
measure, not the realness of the zeros itself.
\end{remark}

\begin{corollary}[valid forms of the main-direction implication]
\label{v4-arith:w5-j:cor:equivalent-forms}
\ClaimStatusProvedHere. Each of the following hypotheses suffices for RH:
\begin{enumerate}
\item (HB\(^{+}\)), where (HB\(^{+}\)) includes real-part divisor
control over \(\Xi\) (reduces to
\Cref{v4-arith:w5-j:cor:no-UHP-zeros}).
\item The full conjunction (SH) + (Det) + (HB\(^{+}\)) + (BD\(^{+}\)).
\item A de Branges formulation in which the real part of the specified
Hermite--Biehler function has divisor \(\mathrm{div}\,\Xi\).
\item The Conrey--Li 2000 construction: \(\mathcal{H}(E_{\xi})\)
admitting a self-adjoint multiplication operator whose spectral
measure is the zero-counting measure on
\(\mathcal{Z}_{\xi}\cap\mathbb{R}\).
\end{enumerate}
No assertion is made that (SH) + (Det) + (BD\(^{+}\)) alone gives
zero placement without the Hermite--Biehler or real-spectrum input. The
main theorem proves the sufficient implication stated above.
\end{corollary}

\begin{proof}
(1): \Cref{v4-arith:w5-j:thm:zeros-of-Xi-under-HB}. (2): the theorem
itself. (3): Lagarias Acta Arith.~89 (1999), with the de Branges
datum specified. (4): Conrey--Li
arXiv:math/9812166 Thm.~2, with real spectrum part of the hypothesis.
\end{proof}

\begin{theorem}[main direction as a component of
\Cref{v4-arith:platonic:thm:RH-completion}]
\label{v4-arith:w5-j:thm:main-direction-as-component}
\ClaimStatusProvedHere. The positive Hamburger and sampling
spectral hypothesis stated in
\Cref{v4-arith:platonic:def:RH-package} and used in
\Cref{v4-arith:platonic:thm:RH-completion} is precisely the
conjunction (SH) + (Det) + (HB\(^{+}\)) + (BD\(^{+}\)) of
\Cref{v4-arith:w5-j:def:spectral-package}.
\Cref{v4-arith:w5-j:thm:VBstar-main-direction} records the implication
from this conjunction to RH inside the P1, P2, P3 framework of
\Cref{v4-arith:platonic:thm:RH-completion}.
\end{theorem}

\begin{proof}
The spectral hypothesis of
\Cref{v4-arith:platonic:def:RH-package} names ``the positive
Hamburger spectral condition: VB\(^{*}\) on the primary
Fock, equivalently the strengthened HB\(^{+}\)/BD\(^{+}\) forms of
\Cref{v4-arith:hpvbconv:thm:VBstar-converse-spectral-regularity}'';
this unpacks to (SH), (Det), (HB\(^{+}\)), (BD\(^{+}\)) by
\Cref{v4-arith:hpvbconv:cor:three-strictly-finer}. The implication is
\Cref{v4-arith:w5-j:thm:VBstar-main-direction}.
\end{proof}

\section{Analytic Hypotheses on \(\xi_{\mathbb R}\)}
\label{v4-arith:w5-j:what-remains}

\begin{remark}[analytic content]
\label{v4-arith:w5-j:rem:open-core}
\Cref{v4-arith:w5-j:thm:VBstar-main-direction} proves the main-direction
implication from the four hypotheses. It does \emph{not} prove RH, because
none of HB\(^{+}\), BD\(^{+}\), or (Det) is proved for \(\xi_{\mathbb
R}\) by the Hadamard product alone. The three analytic assertions have
different logical strength:
\begin{itemize}
\item (HB\(^{+}\)) proof is RH
(\Cref{v4-arith:hpvbconv:thm:de-branges-1994}; de Branges
1994).
\item (BD\(^{+}\)) proof is a quantitative spectral-regularity
input with the Petersson constant fixed
(\Cref{v4-arith:hpvbconv:prop:VBstar-BD-rate}).
\item (Det) proof is the positive regularised moment
functional with \(M_{2n}^{1/(2n)}=O(n)\), the residual
\((\mathrm{R}_{\mathrm{analytic}})\) of
\Cref{v4-arith:hpar-det:thm:single-analytic-residual}.
\end{itemize}
\end{remark}

\begin{theorem}[sufficient residual form of RH]
\label{v4-arith:w5-j:thm:residual-form}
\ClaimStatusProvedHere. Under P2, P3, H-P\(^{\mathrm{Ar}}\), the
conjunction of the four conditions (SH), (Det), (HB\(^{+}\)), (BD\(^{+}\))
on the primary Fock implies the Riemann Hypothesis. The converse is not
asserted here.
\end{theorem}

\begin{proof}
This is exactly
\Cref{v4-arith:w5-j:thm:VBstar-main-direction}. The reverse implication
would require the spectral-regularity hypotheses of
\Cref{v4-arith:hpvbconv:thm:VBstar-converse-spectral-regularity}, and
is not part of this theorem.
\end{proof}

\begin{remark}[the analytic arrow]
\label{v4-arith:w5-j:rem:contribution}
The analytic content is the main-direction arrow
\[
(\mathrm{SH})+(\mathrm{Det})+(\mathrm{HB}^{+})+(\mathrm{BD}^{+})\;\Longrightarrow\;\mathrm{RH},
\]
isolating the analytic content as the proof of these four
hypotheses on \(\xi_{\mathbb R}\). The main-direction argument
makes the sufficient chain explicit and identifies the exact shape of
the required analytic theorem.
\end{remark}

\begin{remark}[relation to the structural equivalence]
\label{v4-arith:w5-j:rem:construction2-relation}
\Cref{v4-arith:hpvbconv:thm:VBstar-converse-spectral-regularity}
gives a conditional comparison between Reg, HB\(^{+}\), BD\(^{+}\),
and the VB\(^{*}\) converse under P2, P3, H-P\(^{\mathrm{Ar}}\), and the
stated Fock-space identifications.  The complementary arrow is
(SH) + (Det) + (HB\(^{+}\)) + (BD\(^{+}\)) \(\Rightarrow\) RH.  Together
they say that regularity is the needed input and that the full
spectral hypothesis is sufficient for RH.
\end{remark}

\section{Bibliography Anchors}
\label{v4-arith:w5-j:bibliography}

The primary sources for the main-direction implication, beyond the
operator-theoretic list of \Cref{v4-arith:hpvbconv:sources}, are:

\begin{description}
\item[de Branges, \emph{Hilbert Spaces of Entire Functions}
(Prentice-Hall 1968).] Theorems 22--23 (Hermite--Biehler
structure), theorem 32 (multiplication operator on
\(\mathcal{H}(E)\)). The unconditional theorem that
\(E\) Hermite--Biehler \(\Leftrightarrow\) de Branges kernel
positive-definite, with real separated zeros for the real and
imaginary parts, is the backbone of the main-direction implication.
\item[Lagarias, Acta Arith.~89 (1999), 217--234.] Reformulation:
RH \(\Leftrightarrow\) the Lagarias function \(E_{\xi}\) is
Hermite--Biehler. This is (HB\(^{+}\)) \(\Leftrightarrow\) RH, the
single-hypothesis equivalent form used in
\Cref{v4-arith:w5-j:rem:HBplus-alone}.
\item[Lagarias, Acta Arith.~89 (1999).] Survey of de Branges
reformulations; \S4.2 on the exponential type of \(E_{\xi}\); \S5
on the Beurling--Malliavin density of \(\mathcal{Z}_{\xi}\); \S6
on the Conrey--Li counter-example discussion.
\item[Conrey--Li, arXiv:math/9812166 (2000).] Theorem~2: the
de Branges space \(\mathcal{H}(E_{\xi})\) does not in general
admit an arithmetic multiplication operator realising RH; the
authors construct explicit counter-examples for \(L\)-functions
where the Hermite--Biehler condition fails (e.g.~for certain
Dirichlet \(L\)-functions with Siegel-type zeros). This confirms
that (HB\(^{+}\)) is a substantive hypothesis not automatic from
the Hadamard product.
\item[Beurling--Malliavin, Acta Math.~107 (1962), 291--309.]
Density theorem: completeness radius of
\(\{e^{i\lambda_{k}x}\}\) equals Beurling--Malliavin density of
\(\{\lambda_{k}\}\). This is the analytical basis for the
(BD\(^{+}\)) \(\Rightarrow\) sampling implication.
\item[Hamburger, Math.~Ann.~81 (1920), 235--319.] Existence
theorem: positive-semidefinite Hankel \(\Leftrightarrow\) positive
Borel measure on \(\mathbb{R}\). Hamburger's own proof via the
Helly selection theorem and the Banach--Alaoglu compactness of
probability measures on \(\mathbb{R}\).
\item[Carleman, Comptes Rendus Acad.~Sci.~Paris (1922); also in
\emph{Les fonctions quasi analytiques} (Gauthier-Villars 1926).]
Determinacy criterion: \(\sum M_{2n}^{-1/(2n)}=\infty\)
\(\Rightarrow\) Hamburger determinacy. The proof uses
quasi-analytic classes of entire functions and constitutes the
classical foundation for the determinacy condition (Det).
\item[Koosis, \emph{The Logarithmic Integral II} (Cambridge 1992),
Ch.~VIII.] Modern treatment of Beurling--Malliavin density and
its applications to completeness of exponential systems.
\item[Levin, \emph{Distribution of Zeros of Entire Functions}
(AMS 1980), Ch.~V--VII.] Classical reference on Hermite--Biehler
and Laguerre--P\'olya class; the structure theorems used in the
HB\(^{+}\) characterisation.
\end{description}

These sources furnish the classical foundations for the implication.
None supplies the proof of HB\(^{+}\), BD\(^{+}\),
or (Det) on \(\xi_{\mathbb R}\).

\section{Spectral Boundary}
\label{v4-arith:w5-j:summary}

\begin{enumerate}
\item \(\xi_{\mathbb R}\) is entire of order \(1\), Hadamard genus
\(1\), with \(\xi_{\mathbb R}(s)=\xi_{\mathbb R}(1-s)\); hence
\(\Xi(z)=\xi_{\mathbb R}(\tfrac{1}{2}+iz)\) has real Taylor
coefficients and a zero set symmetric under the Klein four group
(\Cref{v4-arith:w5-j:prop:xi-order}, \Cref{v4-arith:w5-j:prop:func-eq}).
\item The regularised moment sequence \(\{M_{n}\}\) of the zero
multi-set \(\mathcal{Z}_{\xi}\) is real, with odd moments vanishing
(\Cref{v4-arith:w5-j:prop:moments-real}).
\item Positive Hamburger moment structure (SH) plus
Hamburger determinacy (Det) give a unique positive Borel measure
\(\nu_{\xi}\) on \(\mathbb{R}\) with moments \(\{M_{n}\}\)
(\Cref{v4-arith:w5-j:cor:unique-measure}).
\item Hermite--Biehler positivity (HB\(^{+}\)), including
real-part divisor control over \(\Xi\), forces
\(\mathcal{Z}_{\xi}\subset\mathbb{R}\)
(\Cref{v4-arith:w5-j:cor:no-UHP-zeros}).
\item Beurling--Malliavin super-critical density (BD\(^{+}\)) gives
the uniqueness-set property for \(\mathcal{H}(E_{\xi})\); a full
sampling formula requires the stronger sampling inequality
(\Cref{v4-arith:w5-j:cor:zeros-sample}).
\item The conjunction (SH) + (Det) + (HB\(^{+}\)) + (BD\(^{+}\))
implies zero-placement and hence RH
(\Cref{v4-arith:w5-j:thm:VBstar-main-direction}).
\item The analytic content is the proof of HB\(^{+}\),
BD\(^{+}\), and (Det) on \(\xi_{\mathbb R}\); these are the spectral
inputs
(\Cref{v4-arith:w5-j:rem:open-core}).
\item The analytic implication matches the RH completion
hypothesis of \Cref{v4-arith:platonic:def:RH-package}: the positive
Hamburger spectral component of that hypothesis is
exactly (SH) + (Det) + (HB\(^{+}\)) + (BD\(^{+}\)), and the arrow to
RH is \Cref{v4-arith:w5-j:thm:VBstar-main-direction}.
\end{enumerate}
